% ============================================================================
% cross_references_demo.tex - Cross-Referencing in LaTeX
% ============================================================================
% Demonstrates the full range of cross-reference commands:
%   \label / \ref         — basic references
%   \autoref              — type-aware references (hyperref)
%   \cref / \Cref         — smart references (cleveref)
%   \nameref              — reference by name/title
%   \pageref              — reference by page number
%
% Targets: sections, equations, figures, tables, theorems, listings
%
% Compilation:
%   pdflatex cross_references_demo.tex
%   pdflatex cross_references_demo.tex   (second pass resolves labels)
% ============================================================================

\documentclass[12pt,a4paper]{article}

% ============================================================================
% PACKAGES
% ============================================================================

\usepackage[utf8]{inputenc}
\usepackage[T1]{fontenc}
\usepackage{lmodern}
\usepackage[english]{babel}
\usepackage[margin=1in]{geometry}
\usepackage{microtype}

% Mathematics
\usepackage{amsmath, amssymb, amsthm}

% Figures and tables
\usepackage{graphicx}
\usepackage{float}
\usepackage{booktabs}

% Code listings (used as a reference target)
\usepackage{listings}
\lstset{basicstyle=\ttfamily\small, frame=single, numbers=left,
        numberstyle=\tiny\color{gray}, breaklines=true}

% Colors
\usepackage{xcolor}

% --- hyperref must be loaded BEFORE cleveref ---
\usepackage[
    colorlinks=true,
    linkcolor=blue,
    citecolor=teal,
    urlcolor=cyan,
    pdftitle={Cross-Referencing in LaTeX},
    pdfauthor={LaTeX Student}
]{hyperref}

% --- cleveref: context-aware cross-references ---
% Must be loaded AFTER hyperref.
\usepackage{cleveref}
% Customise cleveref labels (optional):
\crefname{table}{Table}{Tables}
\crefname{figure}{Figure}{Figures}
\crefname{equation}{Equation}{Equations}
\crefname{section}{Section}{Sections}
\crefname{listing}{Listing}{Listings}

% ============================================================================
% THEOREM ENVIRONMENTS (needed so cleveref knows their type)
% ============================================================================

\newtheorem{theorem}{Theorem}[section]
\newtheorem{lemma}[theorem]{Lemma}

% ============================================================================
% DOCUMENT METADATA
% ============================================================================

\title{\textbf{Cross-Referencing in \LaTeX}\\[0.4em]
       \large \texttt{\textbackslash ref}, \texttt{\textbackslash autoref},
       \texttt{\textbackslash cref}, \texttt{\textbackslash nameref},
       \texttt{\textbackslash pageref}}
\author{LaTeX Student}
\date{\today}

% ============================================================================
% DOCUMENT
% ============================================================================

\begin{document}

\maketitle
\tableofcontents
\newpage

% ============================================================================
% SECTION 1: INTRODUCTION
% ============================================================================

\section{Introduction}
\label{sec:intro}

Cross-referencing ties different parts of a document together.  \LaTeX's
\verb|\label| / \verb|\ref| pair provides the foundation; the
\texttt{hyperref} package turns every reference into a clickable hyperlink
and adds \verb|\autoref|; the \texttt{cleveref} package adds \verb|\cref|
and \verb|\Cref| which automatically prepend the correct type word (Section,
Figure, Table, \ldots).

\medskip
\noindent This document defines one label in each of the following
categories and then demonstrates every reference command against each label:
\begin{enumerate}
    \item A section (\cref{sec:targets})
    \item A numbered equation (\cref{eq:euler})
    \item A figure (\cref{fig:parabola})
    \item A table (\cref{tab:data})
    \item A theorem (\cref{thm:pythagoras})
    \item A code listing (\cref{lst:fibonacci})
\end{enumerate}

% ============================================================================
% SECTION 2: LABEL TARGETS
% ============================================================================

\section{Label Targets}
\label{sec:targets}

This section defines all the \verb|\label| anchors that the later sections
will reference.

\subsection{An Equation}
\label{subsec:eq}

Euler's identity is arguably the most beautiful equation in mathematics:
\begin{equation}
    \label{eq:euler}
    e^{i\pi} + 1 = 0.
\end{equation}

The quadratic formula,
\begin{equation}
    \label{eq:quadratic}
    x = \frac{-b \pm \sqrt{b^2 - 4ac}}{2a},
\end{equation}
solves any polynomial $ax^2 + bx + c = 0$.

An aligned multi-line equation:
\begin{align}
    (a+b)^2 &= a^2 + 2ab + b^2,  \label{eq:expand_sq} \\
    (a-b)^2 &= a^2 - 2ab + b^2.  \label{eq:expand_diff}
\end{align}

\subsection{A Figure}
\label{subsec:fig}

\begin{figure}[h]
\centering
% Inline TikZ parabola (no external file needed)
\begin{minipage}{0.5\textwidth}
\centering
\setlength{\unitlength}{1cm}
\begin{picture}(6,5)
    % Axes
    \put(0,2.5){\vector(1,0){6}}
    \put(3,0){\vector(0,1){5}}
    % Curve: y = (x-3)^2 / 2 + 0.5, x from 0.5 to 5.5
    % Approximated with a Bezier via multiple line segments:
    \multiput(0.6,4.5)(0.3,-0.3){4}{\line(3,-2){0.3}}
    \multiput(1.8,3.3)(0.3,-0.15){3}{\line(2,-1){0.3}}
    \put(2.7,2.85){\line(1,0){0.6}}
    \multiput(3.3,2.85)(0.3,0.15){3}{\line(2,1){0.3}}
    \multiput(4.2,3.3)(0.3,0.3){4}{\line(3,2){0.3}}
    % Labels
    \put(5.8,2.3){\small$x$}
    \put(3.1,4.8){\small$y$}
    \put(3.3,2.7){\small$V$}
\end{picture}
\end{minipage}
\caption{Parabola $y = (x-x_0)^2/(2p) + y_0$ with vertex $V$.}
\label{fig:parabola}
\end{figure}

\subsection{A Table}
\label{subsec:tab}

\begin{table}[h]
\centering
\caption{Summary of cross-reference commands and their packages.}
\label{tab:data}
\begin{tabular}{lll}
\toprule
\textbf{Command} & \textbf{Package} & \textbf{Output example} \\
\midrule
\verb|\ref{label}|        & \textit{kernel}      & 2 \\
\verb|\autoref{label}|    & \texttt{hyperref}    & Section~2 \\
\verb|\cref{label}|       & \texttt{cleveref}    & section~2 \\
\verb|\Cref{label}|       & \texttt{cleveref}    & Section~2 \\
\verb|\nameref{label}|    & \texttt{hyperref}    & Label Targets \\
\verb|\pageref{label}|    & \textit{kernel}      & 2 \\
\bottomrule
\end{tabular}
\end{table}

\subsection{A Theorem}
\label{subsec:thm}

\begin{theorem}[Pythagorean Theorem]
\label{thm:pythagoras}
In any right triangle with legs $a$, $b$ and hypotenuse $c$:
\[
    a^2 + b^2 = c^2.
\]
\end{theorem}

\begin{lemma}
\label{lem:square_sum}
For positive integers $a < b < c$ with $a^2 + b^2 = c^2$, both $a$ and $b$
cannot be odd simultaneously.
\end{lemma}

\subsection{A Code Listing}
\label{subsec:lst}

\begin{lstlisting}[language=Python, caption={Iterative Fibonacci sequence generator.}, label=lst:fibonacci]
def fibonacci(n):
    """Yield the first n Fibonacci numbers."""
    a, b = 0, 1
    for _ in range(n):
        yield a
        a, b = b, a + b

print(list(fibonacci(10)))
# Output: [0, 1, 1, 2, 3, 5, 8, 13, 21, 34]
\end{lstlisting}

% ============================================================================
% SECTION 3: REFERENCE COMMANDS
% ============================================================================

\section{Reference Commands}
\label{sec:refcmds}

\subsection{\texttt{\textbackslash ref} — Basic Reference}
\label{subsec:ref}

\verb|\ref{label}| produces only the number of the referenced item, with no
type word.  You must write the type word yourself.

\begin{itemize}
    \item Section~\ref{sec:targets} contains all label targets.
    \item Equation~(\ref{eq:euler}) is Euler's identity.
    \item Figure~\ref{fig:parabola} shows a parabola.
    \item Table~\ref{tab:data} lists the reference commands.
    \item Theorem~\ref{thm:pythagoras} is the Pythagorean theorem.
    \item Listing~\ref{lst:fibonacci} shows a Fibonacci generator.
\end{itemize}

\noindent\textbf{Best practice:} use a non-breaking space \verb|~| between
the type word and \verb|\ref| to prevent line breaks inside a reference.

\subsection{\texttt{\textbackslash autoref} — Hyperref's Type-Aware Reference}
\label{subsec:autoref}

\verb|\autoref{label}| automatically prepends the type word by inspecting
the context in which \verb|\label| was used.  It uses hyperref's internal
mapping of label prefixes to type strings.

\begin{itemize}
    \item \autoref{sec:targets} (section — uses \texttt{sec:} prefix)
    \item \autoref{eq:euler} (equation — uses \texttt{eq:} prefix)
    \item \autoref{fig:parabola} (figure — uses \texttt{fig:} prefix)
    \item \autoref{tab:data} (table — uses \texttt{tab:} prefix)
    \item \autoref{thm:pythagoras} (theorem — uses \texttt{thm:} prefix)
    \item \autoref{lst:fibonacci} (listing — uses \texttt{lst:} prefix)
\end{itemize}

\noindent\textbf{Limitation:} \verb|\autoref| deduces the type from the
label prefix, so you must use consistent prefixes (\texttt{sec:},
\texttt{eq:}, \texttt{fig:}, \texttt{tab:}, \ldots).  It also cannot
produce \textit{plural} forms automatically.

\subsection{\texttt{\textbackslash cref} and \texttt{\textbackslash Cref} — cleveref}
\label{subsec:cref}

\texttt{cleveref}'s \verb|\cref| is the most powerful option.  It:
\begin{itemize}
    \item Automatically prepends the correct type word (\verb|\cref| in
          lowercase, \verb|\Cref| with a capital first letter)
    \item Handles \textbf{multiple references} in one command, with smart
          range compression
    \item Works regardless of label prefixes
\end{itemize}

\subsubsection{Single References}

\begin{itemize}
    \item \cref{sec:targets} \quad (lowercase, mid-sentence)
    \item \Cref{sec:targets} \quad (capitalised, sentence-initial)
    \item \cref{eq:euler}
    \item \cref{fig:parabola}
    \item \cref{tab:data}
    \item \cref{thm:pythagoras}
    \item \cref{lst:fibonacci}
\end{itemize}

\subsubsection{Multiple References}

\begin{itemize}
    \item Two equations: \cref{eq:euler,eq:quadratic}
    \item Three equations: \cref{eq:euler,eq:quadratic,eq:expand_sq}
    \item A range: \cref{eq:expand_sq,eq:expand_diff}
          (cleveref prints ``equations~\ref{eq:expand_sq}
           and~\ref{eq:expand_diff}'' or a range if consecutive)
    \item Mixed types: \cref{sec:targets,fig:parabola,tab:data}
          (cleveref groups by type automatically)
\end{itemize}

\subsection{\texttt{\textbackslash nameref} — Reference by Title}
\label{subsec:nameref}

\verb|\nameref{label}| (provided by \texttt{hyperref}) prints the
\emph{name} — the section heading, caption, theorem name, etc. — rather
than its number.

\begin{itemize}
    \item \nameref{sec:intro} (section title)
    \item \nameref{sec:targets} (section title)
    \item \nameref{fig:parabola} (figure caption)
    \item \nameref{tab:data} (table caption)
    \item \nameref{thm:pythagoras} (theorem heading)
\end{itemize}

\noindent\verb|\nameref| is particularly useful in introductions when you
want to refer to a section by its full descriptive title.

\subsection{\texttt{\textbackslash pageref} — Reference by Page Number}
\label{subsec:pageref}

\verb|\pageref{label}| prints the page number of the referenced item.
Combine it with \verb|\ref| for ``Number on page N'' style references:

\begin{itemize}
    \item \cref{sec:targets} is on page~\pageref{sec:targets}.
    \item \cref{eq:euler} appears on page~\pageref{eq:euler}.
    \item \cref{fig:parabola} is on page~\pageref{fig:parabola}.
    \item \cref{tab:data} is on page~\pageref{tab:data}.
    \item \cref{thm:pythagoras} is on page~\pageref{thm:pythagoras}.
    \item \cref{lst:fibonacci} begins on page~\pageref{lst:fibonacci}.
\end{itemize}

% ============================================================================
% SECTION 4: BEST PRACTICES
% ============================================================================

\section{Best Practices}
\label{sec:best_practices}

\subsection{Label Naming Conventions}

Consistent prefixes make it easy to see at a glance what kind of object a
label refers to, and they are required for \verb|\autoref| to work correctly:

\begin{center}
\begin{tabular}{ll}
\toprule
\textbf{Object type}  & \textbf{Recommended prefix} \\
\midrule
Part          & \texttt{part:} \\
Chapter       & \texttt{chap:} \\
Section       & \texttt{sec:} \\
Subsection    & \texttt{subsec:} \\
Equation      & \texttt{eq:} \\
Figure        & \texttt{fig:} \\
Table         & \texttt{tab:} \\
Theorem/Lemma & \texttt{thm:} / \texttt{lem:} \\
Algorithm     & \texttt{alg:} \\
Listing       & \texttt{lst:} \\
Appendix      & \texttt{app:} \\
\bottomrule
\end{tabular}
\end{center}

\subsection{Where to Place \texttt{\textbackslash label}}

\begin{itemize}
    \item \textbf{Sections:} immediately after the \verb|\section| command.
    \item \textbf{Equations:} inside the equation environment, before
          \verb|\end{equation}|.
    \item \textbf{Figures/Tables:} immediately after \verb|\caption{}|
          (not before—this is a common source of ``wrong page'' bugs).
    \item \textbf{Theorems:} immediately after the
          \verb|\begin{theorem}[...]| line.
    \item \textbf{Listings:} in the optional argument of \verb|\lstlisting|:
          \verb|[..., label=lst:name]|.
\end{itemize}

\subsection{Choosing the Right Command}

\begin{description}
    \item[\texttt{\textbackslash cref} (cleveref)] Recommended for all new
        documents. Handles singular/plural, ranges, and mixed types
        automatically.
    \item[\texttt{\textbackslash autoref} (hyperref)] Good fallback if
        \texttt{cleveref} is not available.  Requires consistent label
        prefixes.
    \item[\texttt{\textbackslash ref}] Use only when you need the bare number
        and full explicit control (e.g., inside a macro).
    \item[\texttt{\textbackslash nameref}] Use in prose when the full section
        title provides more context than a number.
    \item[\texttt{\textbackslash pageref}] Use in documents where readers
        may print the PDF and need physical page guidance.
\end{description}

% ============================================================================
% SUMMARY
% ============================================================================

\section*{Summary}

This document has demonstrated six cross-reference mechanisms:

\begin{enumerate}
    \item \verb|\label| / \verb|\ref| — kernel commands, bare number output
    \item \verb|\autoref| — hyperref, type-aware, single reference
    \item \verb|\cref| / \verb|\Cref| — cleveref, type-aware, plural and range-smart
    \item \verb|\nameref| — hyperref, prints section/caption title
    \item \verb|\pageref| — kernel, prints page number
\end{enumerate}

All labels were placed on: a section (\cref{sec:targets}), three equations
(\cref{eq:euler,eq:quadratic,eq:expand_sq,eq:expand_diff}), a figure
(\cref{fig:parabola}), a table (\cref{tab:data}), a theorem
(\cref{thm:pythagoras}), and a code listing (\cref{lst:fibonacci}).

\end{document}

% ============================================================================
% COMPILATION NOTES
% ============================================================================
% Two pdflatex passes resolve all cross-references:
%   pdflatex cross_references_demo.tex
%   pdflatex cross_references_demo.tex
%
% latexmk handles this automatically:
%   latexmk -pdf cross_references_demo.tex
% ============================================================================
